\section{Preliminaries}\label{sec:prelim}

\paragraph{Linear programs.} An LP in standard form is $\max\ c^\top x$ s.t.\ $Ax \le b$, $x \ge 0$, with constraint matrix $A\in\R^{m\times n}$ ($m$ constraints, $n$ variables), right-hand side $b\in\R^m$ and objective $c\in\R^n$. Its dual is $\min\ b^\top y$ s.t.\ $A^\top y \ge c$, $y \ge 0$, and when both are feasible their optima are equal (strong duality) \cite{bertsimas1997}. Minimisation, $\ge$ rows, equalities and bounds reduce to this form by negating and adding rows; Lemma~\ref{lem:ops} shows that this preserves TU. $L$ is the number of bits needed to write down $A$, $b$ and~$c$.

\paragraph{Vertices and bases.} A point of the polyhedron $P=\{x : A'x\le b'\}$ (with $x\ge 0$ among the rows) is a \emph{vertex} if it is not the midpoint of two other points of $P$; equivalently, it is the unique solution of $n$ linearly independent constraints tight at it. A feasible, bounded LP over such a $P$ has an optimal vertex \cite{bertsimas1997}. Simplex works in equality form $[A\ \ I]\,\bar x = b$, $\bar x\ge 0$ (with slack variables). A \emph{basis} $B$ is a set of $m$ linearly independent columns of $[A\ \ I]$, $A_B$ is the corresponding square submatrix, and the basic solution sets $x_B = A_B^{-1}b$ and all other variables to~$0$. The basic feasible solutions are exactly the vertices.

\paragraph{Exact versus integral.} The ILP adds $x\in\Z^n$. We call the relaxation \emph{exact} if $\optlp = \optilp$, and an LP solution \emph{integral} if all its coordinates are integers. An integral optimal LP solution proves exactness, but an exact LP can also have fractional optimal solutions. E1 and E3 measure the two properties separately.

\paragraph{Total unimodularity.} $A$ is TU if every square submatrix (any $k$ rows and any $k$ columns) has determinant in $\{-1,0,1\}$. In particular every entry is $-1$, $0$ or $1$, and TU depends on $A$ alone, not on $b$ or $c$. The VC LP's matrix is the \emph{incidence matrix} of the graph: one row per edge, one column per vertex, and a $1$ where the vertex is an endpoint of the edge. (In graph contexts $n=|V|$ and $m=|E|$, which here agrees with the LP convention.) With edges in cyclic order, $A_{C_3}$ has rows $(1,1,0)$, $(0,1,1)$, $(1,0,1)$ and determinant~$2$, so it is not TU. $A_{C_4}$ (Figure~\ref{fig:families}a) has determinant~$0$, and Section~\ref{sec:families} shows that all its square submatrices have determinant in $\{-1,0,1\}$. $C_k$ denotes the cycle on $k$ vertices.
